% !TeX root = ../main.tex

Datenbanken sind ein fester Bestandteil moderner Anwendungen und werden zur Speicherung und Verwaltung von Daten
eingesetzt.
Insbesondere bei Anwendungen, bei denen mehrere Zugriffe gleichzeitig auf Daten erfolgen, müssen Transaktionen
zuverlässig ausgeführt werden, um die Konsistenz der gespeicherten Daten sicherzustellen.
Transaktionen können grundsätzlich trotz korrekter Implementierung aufgrund von temporären Fehlern
fehlschlagen.
Solche Fehler sind beispielsweise Deadlocks, Serialization Failures und Lock Timeouts.
Diese Fehler entstehen häufig durch konkurrierende Zugriffe auf gemeinsame Daten und bedeuten nicht zwangsläufig,
dass eine Transaktion dauerhaft nicht ausführbar ist.
Ein fehlgeschlagener Transaktionsversuch muss daher nicht zwingend zu einem endgültigen Abbruch der Verarbeitung
führen.
Ein Retry-Mechanismus setzt an dieser Stelle an und ermöglicht es, fehlgeschlagene Transaktionen automatisch
erneut auszuführen.
Dies kann wiederum die Zuverlässigkeit einer Anwendung erhöhen, ohne dass ein manueller Eingriff erforderlich ist.
Die erneute Ausführung einer Transaktion ist jedoch nicht ohne Kosten.
Jeder weitere Versuch beansprucht erneut Ressourcen der Datenbank und kann die Ausführungszeit der Transaktion
verlängern.
Gleichzeitig kann ein Retry die Konkurrenz zwischen gleichzeitig ausgeführten Transaktionen weiter erhöhen.
Besonders bei Fehlern, die durch konkurrierende Zugriffe entstehen, ist daher nicht garantiert, dass ein
unmittelbarer Wiederholungsversuch erfolgreich ist.
Ein unkontrollierter Einsatz von Retry-Mechanismen kann somit zwar die Erfolgsrate erhöhen, gleichzeitig aber die
Performance und die Belastung des Systems negativ beeinflussen.

Es werden die drei Fehlerklassen Deadlocks, Serialization Failures und Lock Timeouts auf einer PostgreSQL Datenbank
erzeugt und wiederholt.
Für die Untersuchung werden verschiedene Retry-Strategien mit variierenden Retry-Counts und Delays eingesetzt und
bei unterschiedlichen Concurrency-Werten getestet.
Bewertet werden die Strategien, insbesondere anhand der Erfolgsrate und Ausführungszeit der Transaktionen.
Dabei gilt die Erfolgsrate als Maß für die Zuverlässigkeit und die Ausführungszeit für Performance.
Ziel ist es, zu untersuchen wie sich die einzelnen Parameter auswirken und ob sich Unterschiede zwischen den
Fehlerklassen ergeben.
Des Weiteren soll sich ergeben, welche Konfigurationen unter den untersuchten Bedingungen sinnvoll sind und welcher
Kompromisse zwischen Erfolgsrate und Performance entsteht.

Diese Arbeit soll die Forschungsfrage: \glqq Welchen Einfluss haben unterschiedliche Retry-Strategien und deren
Parameter auf die Performance und Zuverlässigkeit von Datenbanken bei definierten Fehlerklassen?\grqq \ beantworten.
Ausgehend von der Problemanalyse werden Anforderungen an einen
Retry-Mechanismus abgeleitet.
Auf dieser Grundlage werden die Experimente konzipiert und in einer kontrollierten Testumgebung durchgeführt.
Die untersuchten Strategien und Parameter werden dabei systematisch variiert.
Anschließend werden die Ergebnisse hinsichtlich der Erfolgsrate und Ausführungszeit ausgewertet und die untersuchten
Konfigurationen bewertet.
Unterstützend werden weitere Messgrößen wie Retry-Overhead und Durchsatz betrachtet.

Die Arbeit ist wie folgt aufgebaut.
Zunächst werden die Grundlagen erklärt und die Arbeit in den Stand der Forschung eingeordnet.
Anschließend wird das Problem analysiert und das Konzept der Untersuchung vorgestellt.
Darauf folgt die Beschreibung der Implementierung und der Durchführung der Experimente.
Danach werden die erzielten Ergebnisse dargestellt und diskutiert.
Abschließend wird die Forschungsfrage beantwortet, die Grenzen der Untersuchung betrachtet und mögliche
Erweiterungen für zukünftige Arbeiten
aufgezeigt.